\section*{Chapitre \ref{ch:g8:the-atmosphere} --- L'atmosphère}

\begin{solution}{exo:g8:the-atmosphere:1}
(a) \omterm{def:g8:the-atmosphere:layers}{stratosphère}~; (b) \omterm{def:g8:the-atmosphere:layers}{troposphère}~; (c) \omterm{def:g8:the-atmosphere:layers}{mésosphère}~; (d) \omterm{def:g8:the-atmosphere:layers}{thermosphère}.
\end{solution}

\begin{solution}{exo:g8:the-atmosphere:2}
(a) \qty{15}{\degreeCelsius}. (b) Environ \qty{-56}{\degreeCelsius}.
(c) Environ \qty{-3}{\degreeCelsius}. (d) Environ \qty{-86}{\degreeCelsius}.
(e) De \qty{11}{km} à environ \qty{20}{km} (et sur quelques kilomètres
vers \qty{50}{km} et juste au-dessus de \qty{86}{km}).
\end{solution}

\begin{solution}{exo:g8:the-atmosphere:3}
(a) Faux~: elle baisse dans la \omterm{def:g8:the-atmosphere:layers}{troposphère}, monte dans la \omterm{def:g8:the-atmosphere:layers}{stratosphère},
baisse dans la \omterm{def:g8:the-atmosphere:layers}{mésosphère} et monte de nouveau dans la \omterm{def:g8:the-atmosphere:layers}{thermosphère}.
(b) Vrai. (c) Faux~: elle se trouve dans la \omterm{def:g8:the-atmosphere:layers}{stratosphère}. (d) Vrai.
\end{solution}

\begin{solution}{exo:g8:the-atmosphere:4}
(a) La \omterm{def:g8:the-atmosphere:layers}{troposphère}. (b) La \omterm{def:g8:the-atmosphere:layers}{stratosphère} (dans l'atmosphère standard, la
\omterm{def:g8:the-atmosphere:layers}{tropopause} est à \qty{11}{km}). (c) La \omterm{def:g8:the-atmosphere:layers}{stratosphère}. (d) La \omterm{def:g8:the-atmosphere:layers}{mésosphère}.
\end{solution}

\begin{solution}{exo:g8:the-atmosphere:5}
(a) $20 - 6.5 \times 2 = 20 - 13 = \qty{7}{\degreeCelsius}$. (b) $20 -
6.5 \times 3 = 20 - 19.5 = \qty{0.5}{\degreeCelsius}$. (c) L'air doit se
refroidir de 20~degrés~: $20 \div 6.5 \approx 3.1$, donc vers
\qty{3.1}{km}.
\end{solution}

\begin{solution}{exo:g8:the-atmosphere:6}
(a) $3 - 0.5 = \qty{2.5}{km}$. (b) $6.5 \times 2.5 = 16.25$~: environ
16~degrés. (c) $18 - 16.25 \approx \qty{2}{\degreeCelsius}$.
\end{solution}

\begin{solution}{exo:g8:the-atmosphere:7}
(a) \qty{500}{hPa}, \qty{250}{hPa} et \qty{125}{hPa}. (b) Près du sol, la
règle est bonne (\qty{500}{hPa} contre \qty{506}{hPa}). Plus haut, elle
donne trop~: \qty{250}{hPa} contre \qty{227}{hPa}, et \qty{125}{hPa}
contre \qty{96}{hPa}. La règle de la division par deux n'est qu'une
indication grossière.
\end{solution}

\begin{solution}{exo:g8:the-atmosphere:8}
(a) $10~\text{hPa} = 10^1$~hPa~; $10^3 \div 10^1 = 10^2$~: cent fois plus
petite. (b) $0.01~\text{hPa} = 10^{-2}$~hPa. (c) $10^3 \div 10^{-2} =
10^5$~: cent mille fois plus petite.
\end{solution}

\begin{solution}{exo:g8:the-atmosphere:9}
(a) $6 \div 1.5 = 4$ fois plus large. (b) En montant, l'air qui l'entoure
appuie de moins en moins~: le gaz à l'intérieur se dilate et le ballon
grossit jusqu'à ce que son enveloppe se déchire. (c) Après l'éclatement du
ballon, le parachute ralentit la chute de la radiosonde, pour qu'elle ne
blesse personne et n'abîme rien.
\end{solution}

\begin{solution}{exo:g8:the-atmosphere:10}
(a) $20 \div 7 \approx 2.9$ fois plus épaisse. (b) Plus proche de celle
des \omterm{def:g3:continents-and-oceans:poles}{pôles}~: \qty{11}{km} est à \qty{4}{km} au-dessus de \qty{7}{km}, mais à \qty{9}{km} au-dessous de \qty{20}{km} (elle décrit les \omterm{def:g6:latitude-longitude:latitude}{latitudes} moyennes, pas l'\omterm{def:g3:continents-and-oceans:equator}{équateur}).
(c) Au départ, la \omterm{def:g8:the-atmosphere:layers}{tropopause} est vers \qty{7}{km}, au-dessous de
l'avion~: l'avion est dans la \omterm{def:g8:the-atmosphere:layers}{stratosphère}. À l'arrivée, la \omterm{def:g8:the-atmosphere:layers}{tropopause} est
vers \qty{20}{km}~: l'avion est dans la \omterm{def:g8:the-atmosphere:layers}{troposphère}.
\end{solution}

\begin{solution}{exo:g8:the-atmosphere:11}
(a) $1 - 506 \div 1013 \approx 0.50$~: \textbf{50~\%}. (b) $1 - 227 \div
1013 \approx 0.78$~: \textbf{78~\%}. (c) $1 - 12 \div 1013 \approx 0.99$~:
\textbf{99~\%}. (d) Environ 1~\%.
\end{solution}

\begin{solution}{exo:g8:the-atmosphere:12}
(a) $6371 \div 15 \approx 425$~km par centimètre. (b) $100 \div 425
\approx 0.24$~cm, soit environ \qty{2.4}{mm}. (c) $11 \div 425 \approx
0.026$~cm, soit environ \qty{0.26}{mm}. (d) $0.26 \div 0.1 \approx 2.6$~:
deux ou trois feuilles de papier.
\end{solution}

\begin{solution}{pb:g8:the-atmosphere:1}
\textbf{1.} $300 \times 10 = 3000$~m.

\textbf{2.} $6000 \div 300 = 20$~minutes~; $\num{30000} \div 300 = 100$
minutes.

\textbf{3.} 100~minutes, soit 1~heure et 40~minutes.

\textbf{4.} À \qty{8}{km}~: la \omterm{def:g8:the-atmosphere:layers}{troposphère} (la température baisse
encore). À \qty{30}{km}~: la \omterm{def:g8:the-atmosphere:layers}{stratosphère}.

\textbf{5.} $15 - (-50) = 65$~degrés.

\textbf{6.} $65 \div 10 = \qty{6.5}{\degreeCelsius}$ par kilomètre.

\textbf{7.} De \qty{12}{km} à \qty{20}{km}, elle reste à
\qty{-56}{\degreeCelsius}. Le refroidissement s'arrête entre \qty{10}{km}
et \qty{12}{km}~: la \omterm{def:g8:the-atmosphere:layers}{tropopause} est vers \qty{11}{km}.

\textbf{8.} Elle monte~: \qty{-55}{\degreeCelsius} à \qty{22}{km},
\qty{-47}{\degreeCelsius} à \qty{30}{km}. Le ballon est dans la
\omterm{def:g8:the-atmosphere:layers}{stratosphère}, réchauffée par l'ozone qui absorbe la lumière ultraviolette
du Soleil.

\textbf{9.} $1013 \div 2 \approx \qty{507}{hPa}$~; elle est atteinte entre
\qty{4}{km} (\qty{617}{hPa}) et \qty{6}{km} (\qty{472}{hPa}).

\textbf{10.} $1013 \div 4 \approx \qty{253}{hPa}$.

\textbf{11.} Entre \qty{10}{km} (\qty{265}{hPa}) et \qty{12}{km}
(\qty{194}{hPa}).

\textbf{12.} À partir de \qty{10}{km}, la pression doit encore baisser de $265 - 253 =
12$~hPa, sur les 71~hPa qu'elle perd en 2~km~: $12 \div 71 \times 2
\approx 0.34$~km. La pression vaut le quart de sa valeur au \omterm{def:g4:world-physical-map:altitude}{niveau de la
mer} vers \textbf{\qty{10.3}{km}}.
\end{solution}
